\section{Analisis Ekonofisika dan Teori Permainan}
Proses ekstraksi parameter interaksi dalam penelitian ini dimulai dengan mengaplikasikan metode \textit{sliding window} berukuran 126 hari bursa untuk menangkap dinamika korelasi antar aset yang bersifat temporal dan dinamis. Ukuran jendela waktu tersebut dipilih secara sengaja untuk memastikan bahwa parameter model yang dihasilkan selalu relevan dengan kondisi pasar terbaru, sekaligus memberikan data historis yang cukup untuk analisis statistik yang valid. Pada setiap jendela waktu, dilakukan perhitungan probabilitas gabungan $P(s_i, s_j)$ untuk setiap pasangan aset guna mengukur derajat ketergantungan statistik serta kecenderungan pergerakan harga secara bersamaan antar emiten. Probabilitas ini kemudian menjadi dasar fundamental bagi perhitungan \textit{Quantum Mutual Information} (QMI), yang secara fisik direpresentasikan sebagai kekuatan interaksi atau konstanta kopling magnetik ($J_{ij}$) antar aset di dalam kerangka Hamiltonian. Melalui pendekatan ekonofisika ini, korelasi finansial yang kompleks dapat direduksi menjadi matriks interaksi yang simetris dan siap untuk dipetakan ke dalam sistem \textit{spin} Ising guna keperluan optimasi kuantum selanjutnya.

Setiap pasangan aset dalam portofolio dianalisis secara mendalam melalui konstruksi matriks \textit{payoff} berukuran $2 \times 2$ yang mengevaluasi rata-rata imbal hasil historis pada berbagai kombinasi keputusan strategis. Komponen nilai yang terkandung di dalam matriks \textit{payoff} ini kemudian digunakan untuk mengklasifikasikan tipe permainan yang terbentuk, seperti permainan koordinasi, anti-koordinasi, atau kategori permainan campuran lainnya. Parameter bias individu ($h_i$) dihitung dengan tingkat presisi tinggi berdasarkan selisih ekspektasi \textit{payoff} marginal yang diperoleh oleh suatu aset saat berada dalam posisi aktif dibandingkan dengan posisi pasif terhadap pasar. Penentuan parameter $J$ dan $h$ yang akurat menjamin bahwa fungsi potensial di dalam kerangka kerja \textit{Exact Potential Game} (EPG) senantiasa selaras dengan objektif maksimisasi utilitas yang diinginkan oleh investor. Melalui integrasi yang harmonis antara interaksi antar pasangan ($J$) dan dorongan bias individu ($h$), model portofolio telah siap untuk ditransformasikan secara formal ke dalam bentuk Hamiltonian Ising yang utuh dan kompatibel dengan algoritma kuantum.
